
\documentclass{article}
\usepackage{colortbl}
\usepackage{makecell}
\usepackage{multirow}
\usepackage{supertabular}

\begin{document}

\newcounter{utterance}

\centering \large Interaction Transcript for game `cladder', experiment `full\_v1.5\_default', episode 5310 with Qwen3.5{-}9B.
\vspace{24pt}

{ \footnotesize  \setcounter{utterance}{1}
\setlength{\tabcolsep}{0pt}
\begin{supertabular}{c@{$\;$}|p{.15\linewidth}@{}p{.15\linewidth}p{.15\linewidth}p{.15\linewidth}p{.15\linewidth}p{.15\linewidth}}
    \# & \multicolumn{2}{c}{Player} && \multicolumn{2}{c}{Game Master} \\
    \hline

    \theutterance \stepcounter{utterance}  
    & & & \multicolumn{4}{p{0.6\linewidth}}{
        \cellcolor[rgb]{0.9,0.9,0.9}{
            \makecell[{{p{\linewidth}}}]{
                \texttt{\tiny{[P1$\langle$GM]}}
                \texttt{You are an expert in causal inference. The following question is not a typical commonsense query, but rather a meticulously designed question created by a professor specializing in causal inference, intended to assess the students' mastery of the course content.} \\
\texttt{Imagine a self{-}contained, hypothetical world with only the following conditions, and without any unmentioned factors or causal relationships: Liking spicy food has a direct effect on occupation and salary. Skill has a direct effect on occupation and salary. Occupation has a direct effect on salary. Skill is unobserved. Method 1: We look at how liking spicy food correlates with salary case by case according to occupation. Method 2: We look directly at how liking spicy food correlates with salary in general. To understand how liking spicy food affects salary, is it more correct to use the Method 1 than Method 2?} \\
\texttt{Start your answer with 'Yes' or 'No', followed by additional reasoning or evidence to support your explanation.} \\
            }
        }
    }
     \\ \\

    \theutterance \stepcounter{utterance}  
    & \multicolumn{4}{p{0.6\linewidth}}{
        \cellcolor[rgb]{0.561,0.737,0.561}{
            \makecell[{{p{\linewidth}}}]{
                \texttt{\tiny{[P1$\rangle$GM]}}
                \texttt{No} \\
\\ 
\texttt{**Reasoning:**} \\
\\ 
\texttt{In this hypothetical world, the causal structure can be represented as follows:} \\
\texttt{1.  **Liking Spicy Food** \$\textbackslash{}rightarrow\$ **Occupation**} \\
\texttt{2.  **Liking Spicy Food** \$\textbackslash{}rightarrow\$ **Salary** (Direct effect)} \\
\texttt{3.  **Skill** \$\textbackslash{}rightarrow\$ **Occupation**} \\
\texttt{4.  **Skill** \$\textbackslash{}rightarrow\$ **Salary**} \\
\texttt{5.  **Occupation** \$\textbackslash{}rightarrow\$ **Salary**} \\
\\ 
\texttt{Here, **Skill** is an unobserved confounder. It influences both **Occupation** and **Salary**, creating a backdoor path between **Liking Spicy Food** and **Salary** if we do not account for the specific pathways correctly. More critically, **Occupation** acts as a mediator in the relationship between Liking Spicy Food and Salary (Spicy Food \$\textbackslash{}rightarrow\$ Occupation \$\textbackslash{}rightarrow\$ Salary), but it is also influenced by the confounder Skill.} \\
\\ 
\texttt{**Analysis of Method 2 (Total Effect):**} \\
\texttt{Method 2 looks at the raw correlation between Liking Spicy Food and Salary. This captures the **total effect**, which includes:} \\
\texttt{*   The direct effect of Liking Spicy Food on Salary.} \\
\texttt{*   The indirect effect via Occupation (Spicy Food \$\textbackslash{}rightarrow\$ Occupation \$\textbackslash{}rightarrow\$ Salary).} \\
\texttt{*   The spurious correlation induced by the confounder **Skill** (since Skill affects both Occupation and Salary, and Occupation is correlated with Spicy Food).} \\
\\ 
\texttt{If the goal is to understand "how liking spicy food affects salary" in the sense of the **causal effect** (specifically the direct effect or the total effect free of confounding bias), Method 2 is biased because it does not control for the confounding variable **Skill**. The observed correlation will be inflated or distorted by the fact that people with high skill (who are unobserved) tend to have high salaries and specific occupations, which might correlate with liking spicy food.} \\
\\ 
\texttt{**Analysis of Method 1 (Stratification by Occupation):**} \\
\texttt{Method 1 stratifies by Occupation. By conditioning on Occupation, we block the path through the mediator. However, this introduces a new problem: **Collider Bias** (or Selection Bias).} \\
\texttt{*   **Occupation** is a collider because it is caused by both **Liking Spicy Food** and **Skill**.} \\
\texttt{*   When we condition on a collider (Occupation), we open a non{-}causal path between its parents: **Liking Spicy Food** \$\textbackslash{}leftrightarrow\$ **Skill**.} \\
\texttt{*   Once this path is opened, the estimated effect of Liking Spicy Food on Salary within each occupation stratum will be biased by the induced correlation with Skill. We are essentially comparing people with the same occupation but different levels of unobserved skill, and the act of conditioning on occupation makes the unobserved skill levels dependent on the spicy food preference.} \\
\\ 
\texttt{**} \\
            }
        }
    }
    & & \\ \\

    \theutterance \stepcounter{utterance}  
    & & & \multicolumn{2}{p{0.3\linewidth}}{
        \cellcolor[rgb]{0.9,0.9,0.9}{
            \makecell[{{p{\linewidth}}}]{
                \texttt{\tiny{[GM$|$GM]}}
                \texttt{No} \\
            }
        }
    }
    & & \\ \\

    \theutterance \stepcounter{utterance}  
    & & & \multicolumn{2}{p{0.3\linewidth}}{
        \cellcolor[rgb]{0.9,0.9,0.9}{
            \makecell[{{p{\linewidth}}}]{
                \texttt{\tiny{[GM$|$GM]}}
                \texttt{no} \\
            }
        }
    }
    & & \\ \\

    \theutterance \stepcounter{utterance}  
    & & & \multicolumn{2}{p{0.3\linewidth}}{
        \cellcolor[rgb]{0.9,0.9,0.9}{
            \makecell[{{p{\linewidth}}}]{
                \texttt{\tiny{[GM$|$GM]}}
                \texttt{game\_result = WIN} \\
            }
        }
    }
    & & \\ \\

\end{supertabular}
}

\end{document}
